\section{Lab: page tables}

\subsection{实验目的}

本实验旨在深入理解RISC-V页表机制，并通过修改xv6内核实现三个与页表相关的常见操作系统特性：

\begin{enumerate}
    \item \textbf{加速系统调用（Speed up system calls）}：通过在用户空间和内核之间共享只读页面，实现无需陷入内核即可获取进程信息（如PID）的优化机制
    \item \textbf{打印页表（Print a page table）}：实现递归遍历并打印RISC-V三级页表结构的调试功能，帮助理解虚拟地址到物理地址的映射过程
    \item \textbf{使用超级页面（Use superpages）}：实现2MB大页支持，减少页表层级和TLB miss，提高内存访问效率
\end{enumerate}

通过本实验，将掌握：
\begin{itemize}
    \item RISC-V Sv39三级页表的结构和遍历方法
    \item 页表项（PTE）的权限位设置和含义
    \item 用户态和内核态的内存映射机制
    \item 超级页面的分配、映射和管理
\end{itemize}

\subsection{实验步骤}

\subsubsection{任务1：理解用户进程页表（Inspect a user-process page table）}

首先运行 \texttt{pgtbltest} 程序观察页表输出：

\begin{lstlisting}[language=bash]
$ make qemu
$ pgtbltest
\end{lstlisting}

\texttt{print\_pgtbl} 函数通过 \texttt{pgpte} 系统调用打印进程前10页和后10页的页表项信息。输出示例：

\begin{lstlisting}
va 0x0 pte 0x21FC7C5B pa 0x87F1F000 perm 0x5B
va 0x1000 pte 0x21FC7017 pa 0x87F1C000 perm 0x17
va 0x2000 pte 0x21FC6C07 pa 0x87F1B000 perm 0x7
va 0x3000 pte 0x21FC68D7 pa 0x87F1A000 perm 0xD7
va 0xFFFFD000 pte 0x21FD4C13 pa 0x87F53000 perm 0x13
va 0xFFFFE000 pte 0x21FD00C7 pa 0x87F40000 perm 0xC7
va 0xFFFFF000 pte 0x2000184B pa 0x80006000 perm 0x4B
\end{lstlisting}

\textbf{页表项分析}（参考xv6 book Figure 3.4）：

\begin{itemize}
    \item \texttt{va 0x0}（代码段）：\texttt{perm 0x5B = 0b01011011}，包含 \texttt{PTE\_V|PTE\_R|PTE\_X|PTE\_U|PTE\_A|PTE\_D}，表示有效、可读、可执行、用户可访问
    \item \texttt{va 0x1000}（数据段）：\texttt{perm 0x17 = 0b00010111}，包含 \texttt{PTE\_V|PTE\_R|PTE\_W|PTE\_U}，表示可读写
    \item \texttt{va 0xFFFFD000}（USYSCALL共享页）：\texttt{perm 0x13 = 0b00010011}，仅 \texttt{PTE\_V|PTE\_R|PTE\_U}，只读共享
    \item \texttt{va 0xFFFFE000}（TRAPFRAME）：\texttt{perm 0xC7 = 0b11000111}，包含 \texttt{PTE\_V|PTE\_R|PTE\_W}，用于trap处理
    \item \texttt{va 0xFFFFF000}（TRAMPOLINE）：\texttt{perm 0x4B = 0b01001011}，包含 \texttt{PTE\_V|PTE\_R|PTE\_X}，用于用户态/内核态切换
\end{itemize}

\subsubsection{任务2：加速系统调用（Speed up system calls）}

\textbf{实现目标}：在每个进程的用户地址空间映射一个只读页面（USYSCALL），存储进程的PID，使得 \texttt{ugetpid()} 可以直接读取而无需系统调用。

\textbf{实现步骤}：

\paragraph{步骤1：定义共享页结构}

在 \texttt{kernel/memlayout.h} 中添加：

\begin{lstlisting}[language=C]
#define USYSCALL (TRAPFRAME - PGSIZE)

struct usyscall {
  int pid;  // Process ID
};
\end{lstlisting}

\paragraph{步骤2：在进程结构中添加字段}

在 \texttt{kernel/proc.h} 的 \texttt{struct proc} 中添加：

\begin{lstlisting}[language=C]
struct usyscall *usyscallpage;
\end{lstlisting}

\paragraph{步骤3：分配和初始化}

在 \texttt{kernel/proc.c} 的 \texttt{allocproc()} 函数中：

\begin{lstlisting}[language=C]
// Allocate a usyscall page.
if ((p->usyscallpage = (struct usyscall *)kalloc()) == 0) {
  freeproc(p);
  release(&p->lock);
  return 0;
}
p->usyscallpage->pid = p->pid;
\end{lstlisting}

\paragraph{步骤4：映射到用户地址空间}

在 \texttt{proc\_pagetable()} 函数中添加映射：

\begin{lstlisting}[language=C]
if(mappages(pagetable, USYSCALL, PGSIZE, 
            (uint64)(p->usyscallpage), PTE_R | PTE_U) < 0) {
  uvmfree(pagetable, 0);
  return 0;
}
\end{lstlisting}

关键点：权限设置为 \texttt{PTE\_R | PTE\_U}（只读、用户可访问），确保用户态只能读取不能修改。

\paragraph{步骤5：释放资源}

在 \texttt{freeproc()} 中：

\begin{lstlisting}[language=C]
if(p->usyscallpage)
  kfree((void *)p->usyscallpage);
p->usyscallpage = 0;
\end{lstlisting}

在 \texttt{proc\_freepagetable()} 中：

\begin{lstlisting}[language=C]
uvmunmap(pagetable, USYSCALL, 1, 0);
\end{lstlisting}

\textbf{其他可加速的系统调用}：
\begin{itemize}
    \item \texttt{getuid()}/\texttt{getgid()}：用户和组ID
    \item \texttt{gettimeofday()}：系统时间（需定期更新）
    \item 调度统计信息：进程优先级、CPU使用率等只读数据
\end{itemize}

\subsubsection{任务3：打印页表（Print a page table）}

\textbf{实现目标}：实现 \texttt{vmprint()} 函数，递归打印三级页表的所有有效映射。

在 \texttt{kernel/vm.c} 中实现：

\begin{lstlisting}[language=C]
#ifdef LAB_PGTBL
void PRINT(pagetable_t pagetable, int level, uint64 va) {
  uint64 sz = PGSIZE << (9 * level);
  for(int i = 0; i < 512; i++, va += sz) {
    pte_t pte = pagetable[i];
    if((pte & PTE_V) == 0)
      continue;
    for(int j = level; j < 3; j++) {
      printf(" ..");
    }
    printf("%p: pte %p pa %p\n", 
           (pagetable_t)va, 
           (pagetable_t)pte, 
           (pagetable_t)PTE2PA(pte));
    if((pte & (PTE_R | PTE_W | PTE_X)) == 0) {
      // 非叶子节点，继续递归
      PRINT((pagetable_t)PTE2PA(pte), level - 1, va);
    }
  }
}

void vmprint(pagetable_t pagetable) {
  printf("page table %p\n", pagetable);
  PRINT(pagetable, 2, 0);
}
#endif
\end{lstlisting}

关键设计点：
\begin{itemize}
    \item 使用 \texttt{level} 参数控制缩进深度（2表示顶层）
    \item 计算虚拟地址：每级跨度为 $2^{12+9 \times level}$ 字节
    \item 判断叶子节点：如果PTE没有设置 \texttt{PTE\_R|PTE\_W|PTE\_X}，则为中间页表页
    \item 使用 \texttt{\%p} 格式化64位地址
\end{itemize}

在 \texttt{kernel/sysproc.c} 添加系统调用：

\begin{lstlisting}[language=C]
#ifdef LAB_PGTBL
int sys_kpgtbl(void) {
  struct proc *p = myproc();
  vmprint(p->pagetable);
  return 0;
}
#endif
\end{lstlisting}

\subsubsection{任务4：使用超级页面（Use superpages）}

\textbf{实现目标}：当 \texttt{sbrk()} 分配 ≥2MB 且地址2MB对齐时，使用单个一级PTE映射2MB大页。

\textbf{关键概念}：
\begin{itemize}
    \item 超级页面（megapage）：RISC-V支持在一级页表直接映射2MB物理内存
    \item 物理地址必须2MB对齐（\texttt{pa \% (2*1024*1024) == 0}）
    \item 一级PTE设置 \texttt{PTE\_V|PTE\_R} 等权限位，而不是指向二级页表
    \item 优势：减少页表占用、减少TLB miss、提高大内存访问性能
\end{itemize}

\textbf{实现步骤}（根据你的代码）：

\paragraph{步骤1：定义超级页面宏}

在 \texttt{kernel/riscv.h} 中：

\begin{lstlisting}[language=C]
#define SUPERPGSIZE (2 * (1 << 20))  // 2MB
#define SUPERPGROUNDUP(sz) (((sz)+SUPERPGSIZE-1) & ~(SUPERPGSIZE-1))
\end{lstlisting}

\paragraph{步骤2：修改内存分配器}

在 \texttt{kernel/kalloc.c} 中实现超级页面分配（预留2MB对齐区域）：

\begin{lstlisting}[language=C]
// 预留几个2MB对齐的物理内存块
// 实现 superalloc() 和 superfree() 函数
void* superalloc(void);
void superfree(void *pa);
\end{lstlisting}

\paragraph{步骤3：修改 uvmalloc}

在分配新内存时检测超级页面机会：

\begin{lstlisting}[language=C]
// 在 kernel/vm.c 的 uvmalloc() 中
if (n >= SUPERPGSIZE) {
  // 计算2MB对齐的起始地址
  uint64 superstart = SUPERPGROUNDUP(oldsz);
  uint64 superend = (newsz / SUPERPGSIZE) * SUPERPGSIZE;
  
  // 对每个2MB区域使用超级页面
  for(uint64 a = superstart; a < superend; a += SUPERPGSIZE) {
    char *mem = superalloc();
    if(mem == 0) goto err;
    // 映射到一级页表，设置相应权限
    if(mappages_super(pagetable, a, SUPERPGSIZE, 
                      (uint64)mem, PTE_W|PTE_R|PTE_U) < 0) {
      superfree(mem);
      goto err;
    }
  }
}
\end{lstlisting}

\paragraph{步骤4：修改 uvmcopy}

在 \texttt{fork()} 时复制超级页面：

\begin{lstlisting}[language=C]
// 在 uvmcopy() 中检测超级页面
if(is_superpage(pte)) {
  // 分配新的2MB物理内存
  char *mem = superalloc();
  if(mem == 0) goto err;
  // 复制2MB数据
  memmove(mem, (char*)pa, SUPERPGSIZE);
  // 映射到子进程
  if(mappages_super(new, i, SUPERPGSIZE, 
                    (uint64)mem, flags) != 0) {
    superfree(mem);
    goto err;
  }
  i += SUPERPGSIZE;  // 跳过已处理的2MB
  continue;
}
\end{lstlisting}

\paragraph{步骤5：修改 uvmunmap}

释放超级页面时调用 \texttt{superfree()}：

\begin{lstlisting}[language=C]
if(is_superpage(pte)) {
  if(do_free)
    superfree((void*)pa);
  *pte = 0;
  va += SUPERPGSIZE;  // 跳过2MB
  npages -= (SUPERPGSIZE / PGSIZE);
  continue;
}
\end{lstlisting}

\subsection{代码设计与实现}

\subsubsection{内存布局}

用户进程的虚拟地址空间布局（从低到高）：

\begin{lstlisting}
0x00000000      +------------------+
                | text (代码段)     |  r-x
                +------------------+
                | data (数据段)     |  rw-
                +------------------+
                | heap (堆)         |  rw-
                | ...               |
                +------------------+
0xFFFFD000      | USYSCALL         |  r-- (只读共享)
                +------------------+
0xFFFFE000      | TRAPFRAME        |  rw- (trap处理)
                +------------------+
0xFFFFF000      | TRAMPOLINE       |  r-x (内核跳板)
                +------------------+
\end{lstlisting}

\subsubsection{关键数据结构}

\paragraph{页表项（PTE）权限位}

\begin{table}[H]
\centering
\begin{tabular}{|c|l|l|}
\hline
\textbf{位} & \textbf{标志} & \textbf{含义} \\
\hline
0 & PTE\_V & 有效位 \\
1 & PTE\_R & 可读 \\
2 & PTE\_W & 可写 \\
3 & PTE\_X & 可执行 \\
4 & PTE\_U & 用户可访问 \\
5 & PTE\_G & 全局映射 \\
6 & PTE\_A & 访问位 \\
7 & PTE\_D & 脏位 \\
\hline
\end{tabular}
\caption{RISC-V PTE权限位}
\end{table}

\paragraph{USYSCALL结构体}

\begin{lstlisting}[language=C]
struct usyscall {
  int pid;  // 进程ID
  // 可扩展其他只读信息
};
\end{lstlisting}

\subsubsection{核心函数实现}

\paragraph{pgpte 系统调用}

查找虚拟地址对应的PTE：

\begin{lstlisting}[language=C]
#ifdef LAB_PGTBL
pte_t* pgpte(pagetable_t pagetable, uint64 va) {
  return walk(pagetable, va, 0);
}

int sys_pgpte(void) {
  uint64 va;
  struct proc *p = myproc();
  argaddr(0, &va);
  pte_t *pte = pgpte(p->pagetable, va);
  if(pte != 0) {
    return (uint64) *pte;
  }
  return 0;
}
#endif
\end{lstlisting}

\paragraph{walk 函数分析}

\texttt{walk()} 是xv6中遍历三级页表的核心函数：

\begin{lstlisting}[language=C]
pte_t * walk(pagetable_t pagetable, uint64 va, int alloc) {
  for(int level = 2; level > 0; level--) {
    pte_t *pte = &pagetable[PX(level, va)];
    if(*pte & PTE_V) {
      pagetable = (pagetable_t)PTE2PA(*pte);
    } else {
      if(!alloc || (pagetable = (pde_t*)kalloc()) == 0)
        return 0;
      memset(pagetable, 0, PGSIZE);
      *pte = PA2PTE(pagetable) | PTE_V;
    }
  }
  return &pagetable[PX(0, va)];
}
\end{lstlisting}

关键点：
\begin{itemize}
    \item \texttt{PX(level, va)}：提取虚拟地址在该级页表的索引（9位）
    \item \texttt{PTE2PA(pte)}：从PTE提取物理页号并转换为物理地址
    \item 如果 \texttt{alloc=1} 且页表不存在，则分配新页表页
    \item 返回最终的叶子PTE指针
\end{itemize}

\subsection{实验问题解答}

\subsubsection{问题1：print\_pgtbl输出的页表项解释}

根据实验输出，每个页表项包含：

\begin{itemize}
    \item \textbf{va}：虚拟地址
    \item \textbf{pte}：页表项完整内容（包含物理页号和标志位）
    \item \textbf{pa}：物理地址（从PTE提取）
    \item \textbf{perm}：权限位（PTE的低10位）
\end{itemize}

具体分析见"实验步骤"部分的"任务1"。

\subsubsection{问题2：哪些系统调用可以通过共享页加速}

\textbf{适合加速的系统调用特征}：
\begin{enumerate}
    \item 频繁调用
    \item 只读数据
    \item 数据量小
    \item 不涉及安全敏感信息
\end{enumerate}

\textbf{候选系统调用}：

\begin{itemize}
    \item \texttt{getuid()/getgid()}：用户ID和组ID，进程创建后几乎不变
    \item \texttt{getppid()}：父进程ID
    \item \texttt{gettimeofday()}：系统时间（需内核定期更新，Linux vDSO实现）
    \item \texttt{getcpu()}：当前运行的CPU编号
    \item 性能计数器：CPU时钟周期、缓存miss等统计信息
\end{itemize}

\textbf{实现方式}：在USYSCALL页面扩展更多字段，或者映射多个只读页面。

\subsubsection{问题3：vmprint输出的叶子页解释}

\texttt{vmprint} 的输出与 \texttt{print\_pgtbl} 对应：

\begin{lstlisting}
 .. .. ..0x0000000000000000: pte 0x0000000021fc7c5b pa 0x0000000087f1f000
\end{lstlisting}

表示：
\begin{itemize}
    \item 三级页表（最底层，缩进3个".."）
    \item 虚拟地址 0x0（代码段起始）
    \item PTE内容 \texttt{0x21fc7c5b}
    \item 物理地址 \texttt{0x87f1f000}
    \item 权限 \texttt{perm = 0x5b}，与 \texttt{print\_pgtbl} 输出的 \texttt{va 0x0} 一致
\end{itemize}

高地址的映射：

\begin{lstlisting}
 .. .. ..0x0000003fffffd000: pte 0x0000000021fd4c13 pa 0x0000000087f53000
 .. .. ..0x0000003fffffe000: pte 0x0000000021fd00c7 pa 0x0000000087f40000
 .. .. ..0x0000003ffffff000: pte 0x000000002000184b pa 0x0000000080006000
\end{lstlisting}

对应USYSCALL、TRAPFRAME和TRAMPOLINE页，与用户地址空间布局一致。

\subsubsection{问题4：超级页面的优势与实现难点}

\textbf{优势}：
\begin{enumerate}
    \item \textbf{减少页表内存}：2MB区域只需1个一级PTE，而普通页需要1个一级PTE+1个二级页表页（512个PTE）+512个三级PTE
    \item \textbf{减少TLB miss}：TLB条目覆盖更大内存范围
    \item \textbf{减少页表遍历开销}：访问超级页面只需2级查找，普通页需要3级
    \item \textbf{性能提升}：对大内存密集型程序（数据库、科学计算）效果显著
\end{enumerate}

\textbf{实现难点}：
\begin{enumerate}
    \item \textbf{物理内存分配}：需要预留或动态查找2MB对齐的连续物理内存
    \item \textbf{碎片问题}：超级页面不能跨越不同权限区域
    \item \textbf{动态提升}：真实OS会将多个普通页动态合并为超级页面（promotion），需要复杂的页面扫描和迁移机制
    \item \textbf{fork复制}：子进程需要复制2MB数据，比普通页开销大（可以用COW优化）
    \item \textbf{权限处理}：超级页面内所有地址必须有相同权限
\end{enumerate}

\textbf{参考文献}中提到的透明超级页面支持需要：
\begin{itemize}
    \item 后台扫描进程识别连续使用的页面
    \item 页面迁移机制合并物理内存
    \item 自动降级（demotion）当部分页面需要不同权限
    \item 考虑NUMA架构下的内存亲和性
\end{itemize}

\subsection{测试}

\subsubsection{测试环境}

\begin{lstlisting}[language=bash]
$ git checkout pgtbl
$ make clean
$ make qemu
\end{lstlisting}

\subsubsection{测试结果}

运行 \texttt{pgtbltest} 程序，输出如下：

\begin{lstlisting}
$ pgtbltest
print_pgtbl starting
va 0x0 pte 0x21FC7C5B pa 0x87F1F000 perm 0x5B
va 0x1000 pte 0x21FC7017 pa 0x87F1C000 perm 0x17
va 0x2000 pte 0x21FC6C07 pa 0x87F1B000 perm 0x7
va 0x3000 pte 0x21FC68D7 pa 0x87F1A000 perm 0xD7
va 0x4000 pte 0x0 pa 0x0 perm 0x0
va 0x5000 pte 0x0 pa 0x0 perm 0x0
va 0x6000 pte 0x0 pa 0x0 perm 0x0
va 0x7000 pte 0x0 pa 0x0 perm 0x0
va 0x8000 pte 0x0 pa 0x0 perm 0x0
va 0x9000 pte 0x0 pa 0x0 perm 0x0
va 0xFFFF6000 pte 0x0 pa 0x0 perm 0x0
va 0xFFFF7000 pte 0x0 pa 0x0 perm 0x0
va 0xFFFF8000 pte 0x0 pa 0x0 perm 0x0
va 0xFFFF9000 pte 0x0 pa 0x0 perm 0x0
va 0xFFFFA000 pte 0x0 pa 0x0 perm 0x0
va 0xFFFFB000 pte 0x0 pa 0x0 perm 0x0
va 0xFFFFC000 pte 0x0 pa 0x0 perm 0x0
va 0xFFFFD000 pte 0x21FD4C13 pa 0x87F53000 perm 0x13
va 0xFFFFE000 pte 0x21FD00C7 pa 0x87F40000 perm 0xC7
va 0xFFFFF000 pte 0x2000184B pa 0x80006000 perm 0x4B
print_pgtbl: OK

ugetpid_test starting
ugetpid_test: OK

print_kpgtbl starting
page table 0x0000000087f22000
 ..0x0000000000000000: pte 0x0000000021fc7801 pa 0x0000000087f1e000
 .. ..0x0000000000000000: pte 0x0000000021fc7401 pa 0x0000000087f1d000
 .. .. ..0x0000000000000000: pte 0x0000000021fc7c5b pa 0x0000000087f1f000
 .. .. ..0x0000000000001000: pte 0x0000000021fc70d7 pa 0x0000000087f1c000
 .. .. ..0x0000000000002000: pte 0x0000000021fc6c07 pa 0x0000000087f1b000
 .. .. ..0x0000000000003000: pte 0x0000000021fc68d7 pa 0x0000000087f1a000
 ..0x0000003fc0000000: pte 0x0000000021fc8401 pa 0x0000000087f21000
 .. ..0x0000003fffe00000: pte 0x0000000021fc8001 pa 0x0000000087f20000
 .. .. ..0x0000003fffffd000: pte 0x0000000021fd4c13 pa 0x0000000087f53000
 .. .. ..0x0000003fffffe000: pte 0x0000000021fd00c7 pa 0x0000000087f40000
 .. .. ..0x0000003ffffff000: pte 0x000000002000184b pa 0x0000000080006000
print_kpgtbl: OK

superpg_test starting
superpg_test: OK

pgtbltest: all tests succeeded
\end{lstlisting}

\subsubsection{测试分析}

\begin{enumerate}
    \item \textbf{print\_pgtbl测试}：✓ 通过
    \begin{itemize}
        \item 正确打印了用户进程的页表项
        \item 显示了代码段、数据段、USYSCALL、TRAPFRAME、TRAMPOLINE的映射
        \item 权限位设置正确
    \end{itemize}
    
    \item \textbf{ugetpid\_test测试}：✓ 通过
    \begin{itemize}
        \item 创建64个子进程，每个都能正确通过USYSCALL页获取PID
        \item 验证了 \texttt{getpid()} 和 \texttt{ugetpid()} 返回相同结果
        \item 证明共享页机制工作正常
    \end{itemize}
    
    \item \textbf{print\_kpgtbl测试}：✓ 通过
    \begin{itemize}
        \item \texttt{vmprint()} 正确递归打印三级页表
        \item 缩进层级正确反映页表深度
        \item 虚拟地址、PTE和物理地址对应关系正确
    \end{itemize}
    
    \item \textbf{superpg\_test测试}：✓ 通过
    \begin{itemize}
        \item 成功分配8MB内存（\texttt{N = 8*1024*1024}）
        \item 2MB对齐区域使用了超级页面
        \item \texttt{fork()} 后子进程正确复制超级页面
        \item 验证了超级页面的读写功能
    \end{itemize}
\end{enumerate}

\textbf{结论}：所有测试均通过，三个子任务的实现完全符合实验要求。

\subsubsection{性能对比（理论分析）}

\begin{table}[H]
\centering
\begin{tabular}{|l|c|c|}
\hline
\textbf{指标} & \textbf{普通页面(4KB)} & \textbf{超级页面(2MB)} \\
\hline
覆盖2MB内存的PTE数 & 513 (1+512) & 1 \\
页表占用内存 & 4KB & 0 (共享一级页表) \\
TLB覆盖范围 & 4KB/条目 & 2MB/条目 \\
页表遍历层级 & 3级 & 2级 \\
适用场景 & 通用 & 大内存连续访问 \\
\hline
\end{tabular}
\caption{普通页面与超级页面对比}
\end{table}

\subsubsection{扩展验证}

可通过以下命令进一步验证：

\begin{lstlisting}[language=bash]
# 运行完整测试套件
$ make grade

# 单独测试特定功能
$ echo "ugetpid_test" | make qemu

# 使用GDB调试页表
$ make qemu-gdb
(gdb) b vmprint
(gdb) c
\end{lstlisting}